\section{Discussion and Open Problems}
\label{sec:discuss}

In this concluding section, we speculate on some possible avenues to attack Average Baby PIH or even PIH itself.

\subsection{Average Baby PIH from Clique Hardness?}

In \cite{Lin21}, the author proved that even constant approximating \kclique{} is $\wone$-hard. In \cite{KK22,CFL+23}, the inapproximability ratio was improved to $k^{o(1)}$. Using CSP words, their results can be described as the following theorem:

\begin{theorem}[\cite{KK22,CFL+23}]
    Assuming $\wone\neq\fpt$, no algorithm can, given a 2CSP instance $\Pi=(X,\Sigma,\Phi)$, distinguish between the following two cases in $f(|X|)\cdot |\Sigma|^{O(1)}$ time, for any computable function $f$:
    \begin{itemize}
        \item $\Pi$ is satisfiable.
        \item The constraints induced by any $|X|^{1-o(1)}$ variables are not simultaneously satisfiable.
    \end{itemize}
\end{theorem}
 Thus, it is natural to ask whether we can get Average Baby PIH by applying direct product construction to a 2CSP instance with the above ``clique-like" soundness? More formally, we have the following open question:

\begin{question}
Is it true that for any integer $r>1$, there is an integer $t>0$, such that for any 2CSP instance $\Pi=(X,\Sigma,\Phi)$, if $\Pi^{\odot t}$ is $r$-average-list satisfiable, then there are $|X|^{1-o(1)}$ variables in $\Pi$ such all constraints amongst them are simultaneously satisfiable?
\end{question}

Note that the counterexample for proving average baby PIH using the direct product construction (Example \ref{ex:1}) does not apply here, since there are $|X|-1$ variables in $\Pi$ that are simultaneously satisfiable.

\subsection{PIH via Direct Product Testing Theorems?}

Our constructed instance in proving Baby PIH is reminiscent of the direct product testing, which has been studied in a recent line of work \cite{DS14, DK17, DFH19}. These results have shown that if the $t$-sized subsets of variables have good local consistency, then there is a global function which agrees with most of the subsets. Formally, we have the following theorem from \cite{DFH19}.

\begin{theorem}
\label{thm:dir}
    There is an absolute constant $C>1$, such that for any $\alpha,\beta \in (0,1)$ with $\alpha+\beta \le 1$, there exists a constant $Q(\alpha,\beta)$, such that given any $k,t,m$ with $k \ge C \cdot t$ and $\alpha t \le m \le (1-\beta) t$ and any finite alphabet $\Sigma$ , we have the following.

    Let $\mathcal F=\{f_S:S \rightarrow \Sigma \mid S \in \binom{[k]}{t}\}$ be an ensemble of functions, one for every size-$t$ subset of $[k]$. Let $\mathcal D(m)$ be the distribution as follows:
    \begin{itemize}
        \item Choose $I \in \binom{[k]}{m}$ uniformly at random.
        \item Choose $A,B$ from the set $\{X \mid X \in \binom{[k]}{t}, X \supseteq I\}$ uniformly at random.
    \end{itemize}

    Suppose the following holds:
    \[
    \Pr_{A,B \sim \mathcal D(m)}[{f_A}_{|A \cap B} = {f_B}_{|A \cap B}] \ge 1-\varepsilon,
    \]
    then there exists a global function $g: [k] \rightarrow \Sigma$ such that
    \[
    \Pr_{A \sim \binom{[k]}{t}}[f_A = g_{|A}] \ge 1-Q(\alpha,\beta) \cdot \varepsilon.
    \]
\end{theorem}

By setting the alphabet of each $f_S$ to be the set of all partial satisfying assignments for $S$ and adding the consistency checks according to the above theorem, one may wonder whether we can ``extract'' a large clique from the $1-Q(\alpha,\beta)\cdot \varepsilon$ fraction of subsets that are globally consistent. Formally, let $\mathcal S=\{A | A \in \binom{[k]}{t} , f_A=g_{|A}\}$ as in \Cref{thm:dir}, can we prove there exists a subset $T \subseteq [k]$ of size $\ge k^{1-o(1)}$, such that the following holds?
\begin{itemize}
    \item For every $(i,j) \in \binom{T}{2}$, there exists $A \in \mathcal S$ with $(i,j) \subseteq A$.
\end{itemize}

However, this might not be true if we only use the size bound on $\mathcal S$ and not the additional structures. Consider the following counter-example:

\begin{example}
    Mark each pair $(i,j) \in \binom{[k]}{2}$ as 1 independently with probability $1-\gamma$ with $\gamma$ to be determined. Take $\mathcal S'$ to be the collection of $t$-sized subsets of $[k]$, with all pairs in it marked as 1: $\mathcal S':=\{A \mid A \in \binom{[k]}{t}, \forall (i,j) \subseteq A, (i,j)\text{ is marked as 1}\}$.

    The probability that a $t$-sized subset belongs to $\mathcal S'$ is $(1-\gamma)^{\binom{t}{2}}$, which can be made arbitrarily close to 1 when we set $\gamma=\gamma(t)$ to be some small enough constant depending on the constant $t$.

    However, it was known that (see e.g. \cite{Grimmett_McDiarmid_1975}) the maximum clique size in this Erd\H{o}s-R\'enyi graph is only $2 \log k/\log (1/(1-\gamma))$, far smaller than $k^{1-o(1)}$.
\end{example}

This example suggests that, in order to potentially prove PIH from direct product testing theorems, one may need to analyze the structure of the collection $\mathcal S$. We leave this as an interesting future direction:

\begin{question}
    Can we prove PIH under $\wone \ne \fpt$ using some appropriate form of direct product testing theorems?
\end{question}